\documentclass[11pt]{article}

\usepackage[T1]{fontenc}
\usepackage[utf8]{inputenc}
\usepackage[a4paper,margin=1in]{geometry}
\usepackage{amsmath,amssymb}
\usepackage{graphicx}
\usepackage{xcolor}
\usepackage[hidelinks]{hyperref}

\title{Draft of Supplementary Discussion 6.7:\\
Graph-Convolutional Newton-Type Reconstruction for PVI}
\author{Hyeuknam Kwon}
\date{}

\newcommand{\R}{\mathbb{R}}
\setcounter{secnumdepth}{4}

\begin{document}

\section*{Main Manuscript}
\begingroup
To improve differential PVI reconstruction, we incorporated a two-stage graph-convolutional Newton-type method (GCNM) \cite{Herzberg2021GCNM}.
At each stage, the nonlinear differential forward solution, voltage residual, Jacobian, and regularized physics-based update were recomputed from the current conductivity estimate.
The current estimate and its corresponding physics-based update were then used as node features on a graph defined over the reconstruction mesh, and a stage-specific GCN produced the conductivity estimate for the next stage.
The two GCN blocks were trained sequentially with independent parameters, so the second stage used physics-based quantities recomputed from the first-stage reconstruction.
The complete formulation and training procedure are described in Supplementary Discussion~6.7.


Because conductivity ground truth was not available for the experimental recordings, we generated a separate synthetic dataset for each ring geometry.
The simulated finger cross-sections contained the principal tissue compartments and two independently varying digital arteries, with variation in finger geometry, tissue conductivity, artery geometry and position, pulse morphology, and measurement conditions.
For every synthetic sample, we evaluated the conductivity field on a fine forward mesh and simulated the corresponding eight-electrode voltages, while using a distinct coarser mesh for the reconstruction target.
Each ring-specific dataset contained 200 virtual anatomies divided into independent training, validation, and test partitions before the beats and samples were generated.
We trained the two graph convolutional networks sequentially using the synthetic conductivity targets and voltages.
After selecting the first-stage network, we used its reconstructions to recompute the nonlinear forward solution, residual, Jacobian, and physics-based update before training the independently parameterized second-stage network.
The complete formulation and training procedure are described in Supplementary Discussion~6.7.
\endgroup

\newpage

\setcounter{section}{6}
\setcounter{subsection}{6}
\setcounter{equation}{16}
\subsection{Graph-convolutional Newton-type reconstruction}
\label{sec:pvi_gcnm}

\begingroup

% Numbered equations continue provisionally from equation (16) in the supplied
% Section 6 text. Supporting equations that are not cited remain unnumbered.

We extended the one-step differential inverse solver described in Supplementary Discussion~6.3 using a graph-convolutional Newton-type method (GCNM) \cite{Herzberg2021GCNM}.
The conventional inverse solver computes one conductivity update using a Jacobian evaluated at the homogeneous baseline.
In contrast, the PVI-GCNM consisted of two reconstruction stages that combined a newly evaluated physics-based update with a learned graph-convolutional mapping.
For the GCNM implementation, we fixed the homogeneous baseline conductivity introduced in Supplementary Discussion~6.3 at $\sigma_0=0.7~\mathrm{S/m}$ for all ring-specific models.

\subsubsection{Stage-wise physics-based update}

We indexed the two reconstruction stages by $n=0,1$ and initialized the differential conductivity as $\Delta\sigma_0=0$.
At each stage, the current conductivity passed to the forward solver was
\begin{equation}
    \sigma_n
    =
    \max\!\left(
        \sigma_0\mathbf{1}+\Delta\sigma_n,\,
        \sigma_{\min}\mathbf{1}
    \right),
    \qquad
    \sigma_{\min}=10^{-4}~\mathrm{S/m}.
\label{eq:gcnm-stage-conductivity}
\end{equation}
Here, the maximum was applied elementwise, and the fixed lower bound $\sigma_{\min}$ prevented zero or negative conductivity values from entering the nonlinear forward solver.
Using the stage conductivity in equation~\eqref{eq:gcnm-stage-conductivity}, the full nonlinear differential forward solution was recomputed with the forward operator $F$ in equation~(11):
\begin{equation*}
    \widehat{\Delta V}_n
    =
    F\!\left(\sigma_n\right)
    -F\!\left(\sigma_0\mathbf{1}\right).
\end{equation*}
The voltage residual was
\begin{equation*}
    r_n
    =
    \widehat{\Delta V}_n
    -\Delta V.
\end{equation*}
The Jacobian $J_n$ was recomputed at $\sigma_n$ rather than being fixed across the two stages.
Using the fine-mesh Laplacian $L$ and the coarse-to-fine mapping $P$ from Supplementary Discussion~6.1, the projected regularization matrix on the reconstruction mesh was
\begin{equation}
    R_{\mathrm{PVI}}
    =
    \frac{1}{2}
    P^{\mathsf T}LP.
\label{eq:gcnm-projected-regularization}
\end{equation}
This is the regularization matrix stored and used by the implemented PVI inverse solver.
It was used directly and was not multiplied by its transpose again.
Using the regularization matrix in equation~\eqref{eq:gcnm-projected-regularization}, the physics-based update $p_n$ for the selected model was obtained by solving
\begin{equation}
    \left(
        J_n^{\mathsf T}J_n
        +\lambda^2R_{\mathrm{PVI}}
    \right)p_n
    =
    -J_n^{\mathsf T}r_n.
\label{eq:gcnm-physics-update}
\end{equation}
We used $\lambda=5\times10^{-4}$.
We did not include a separate Levenberg Marquardt damping term of the form $\lambda_{\mathrm{LM}}I$ because the GCNM physics-based update retained the PVI inverse formulation in equation~(13), which uses Laplacian regularization without additional identity damping.
Thus, $\lambda$ controls the spatial regularization imposed by $R_{\mathrm{PVI}}$ and is not a Levenberg Marquardt damping parameter.
No additional step-size multiplier was applied.

\subsubsection{Graph-convolutional reconstruction}

Each of the $N_K$ triangles in the reconstruction mesh was treated as one graph node.
Two graph nodes were connected when the corresponding triangles shared at least one finite element mesh node.
Let $G$ denote the graph constructed using this edge rule, and let $A_G$ denote its binary adjacency matrix.
The adjacency matrix with self-connections and its corresponding degree matrix were defined by
\begin{equation*}
    \widetilde A_G
    =
    A_G+I,
    \qquad
    \widetilde D_G
    =
    \operatorname{diag}\!\left(\sum_j \widetilde A_{G,ij}\right).
\end{equation*}

Let $c_i\in\R^2$ be the centroid of reconstruction-mesh element $i$, let $c_0$ denote the domain center, and let $\rho$ denote the domain radius.
The dimensionless coordinate features for element $i$ were defined as
\begin{equation*}
    C_i
    =
    \left[
        \bar x_i,\,
        \bar y_i,\,
        \bar r_i
    \right],
    \qquad
    (\bar x_i,\bar y_i)
    =
    \frac{c_i-c_0}{\rho},
    \qquad
    \bar r_i
    =
    \sqrt{\bar x_i^2+\bar y_i^2}.
\end{equation*}
The conductivity scale $s$ was fixed from the $99.5$th percentile of the absolute training targets.
The input features at stage $n$ were the normalized current differential conductivity estimate, the normalized physics-based update, and the coordinate features:
\begin{equation*}
    H_n
    =
    \left[
        \frac{\Delta\sigma_n}{s},\,
        \frac{p_n}{s},\,
        C
    \right].
\end{equation*}

The graph-convolutional layer used in this study was defined as
\begin{equation*}
H_n^{(\ell+1)}
=
g_\ell\!\left(
\widetilde D_G^{-1/2}
\widetilde A_G
\widetilde D_G^{-1/2}
H_n^{(\ell)}
W_n^{(\ell)}
+
\mathbf{1}\bigl(b_n^{(\ell)}\bigr)^{\mathsf T}
\right),
\qquad
H_n^{(0)}=H_n.
\end{equation*}
Here, \(W_n^{(\ell)}\in\mathbb{R}^{F_\ell\times F_{\ell+1}}\) and \(b_n^{(\ell)}\in\mathbb{R}^{F_{\ell+1}}\) are trainable parameters, and \(g_\ell\) is the layer activation.
Thus, \(\mathbf{1}\bigl(b_n^{(\ell)}\bigr)^{\mathsf T}\) applies the same bias vector to every graph node.
Each stage-specific GCN contained three hidden graph-convolutional layers with 64 channels and ReLU activations, followed by a one-channel graph-convolutional layer with linear activation.
Let \(\theta_n\) denote the collection of trainable parameters in the \(n\)-th GCN block.
The two stages used independent parameter sets \(\theta_0\) and \(\theta_1\).
The one-channel output was normalized during training and was multiplied by $s$ to restore conductivity units.
For compact notation, $\operatorname{GCN}_{\theta_n}$ denotes the complete stage-specific mapping, including this rescaling.
The output of each block directly represented the differential conductivity estimate for the next stage:
\begin{equation}
    \Delta\sigma_{n+1}
    =
    \operatorname{GCN}_{\theta_n}\!\left(H_n,G\right),
    \qquad
    n=0,1.
\label{eq:gcnm-stage-output}
\end{equation}
The GCN output in equation~\eqref{eq:gcnm-stage-output} was used directly as the next iterate rather than being added to the current iterate as a residual correction.
The output of the second GCN stage was used as the final reported reconstruction:
\begin{equation*}
    \Delta\sigma_{\mathrm{GCNM}}
    =
    \Delta\sigma_2.
\end{equation*}


\subsubsection{GCNM training}
\label{sec:gcnm-training}

\begingroup

We describe the development of the ring-specific GCNMs in the order in which the training pipeline was constructed and evaluated.
We first explain how the synthetic conductivity and voltage data were generated for each ring geometry and how the training simulations were verified against exact nonlinear forward solutions.
We then describe the training of the two GCNM stages and the controlled comparisons used to select the output formulation, coordinate inputs, and loss coefficients $\alpha$ and $\beta$.
Finally, we define the training loss and the validation metrics used to select the checkpoint for each stage,

\paragraph{Synthetic data generation}
\label{sec:gcnm-synthetic-data}
\mbox{}\par

Direct conductivity ground truth was not available for the experimental PVI recordings.
We therefore generated a separate synthetic training corpus for each ring geometry represented in the study, preserving the principal anatomical and temporal structures expected in a finger cross-section while exposing the reconstruction model to variation in anatomy, tissue conductivity, vascular position, pulse morphology, and measurement conditions.

Each virtual anatomy contained skin, subcutaneous fat, muscle, cortical bone and marrow, ligaments, artery walls, and two blood-filled arteries.
Within this anatomy, the two arteries were represented independently so that their positions, dimensions, orientations, pulse amplitudes, and spatial diffusion into the surrounding muscle could vary between anatomies.
\begingroup
The nominal resting conductivities were taken at $50~\mathrm{kHz}$ from the public IT'IS Tissue Properties Database distributed for Sim4Life \cite{ITIS2024}.
The values were $0.000273~\mathrm{S/m}$ for dry skin, $0.0433~\mathrm{S/m}$ for fat, $0.352~\mathrm{S/m}$ for muscle, $0.0206~\mathrm{S/m}$ for cortical bone, $0.00363~\mathrm{S/m}$ for bone marrow, $0.388~\mathrm{S/m}$ for ligament and tendon, $0.317~\mathrm{S/m}$ for the artery wall, and $0.701~\mathrm{S/m}$ for blood.
These heterogeneous values defined the resting conductivity of each synthetic anatomy and were distinct from the homogeneous $0.7~\mathrm{S/m}$ baseline used by the inverse solver.
\endgroup
Starting from the reference finger model, finger size was varied by up to $15\%$, finger rotation by up to $20^\circ$, and finger position within the ring by up to $1~\mathrm{mm}$ in each spatial direction, while artery size was varied by up to $35\%$, artery position by up to $1.8~\mathrm{mm}$, and artery rotation by up to $30^\circ$.
Tissue conductivities were additionally varied by up to $20\%$, the spatial diffusion surrounding the arteries by up to $75\%$, the pulse-waveform shape by up to $20\%$, and beat duration by up to $15\%$.
These perturbations were drawn independently for each virtual anatomy while retaining the anatomical tissue arrangement and two-artery structure of the reference model.

Let $a$ index a virtual anatomy and $t$ index a temporal sample, and let the superscript $\mathrm f$ identify the fine-mesh conductivity used by the forward solver.
The time-varying conductivity was constructed as
\begin{equation}
    \sigma^{\mathrm f}_{a,t}
    =
    \sigma^{\mathrm f}_{a,\mathrm{rest}}
    +\delta\sigma^{\mathrm f,\mathrm{pulse}}_{a,t}
    +\delta\sigma^{\mathrm f,\mathrm{vascular}}_{a,t}
    +\delta\sigma^{\mathrm f,\mathrm{muscle}}_{a,t}.
\label{eq:gcnm-synthetic-conductivity}
\end{equation}
The pulsatile term in equation~\eqref{eq:gcnm-synthetic-conductivity} described the conductivity change produced by the two arteries over a cardiac cycle.
Seven consecutive 50-sample beats were simulated for each anatomy, with the amplitude of each beat independently multiplied by a factor drawn between 0.70 and 1.30 and with additional perturbations of its morphology and duration.
Across the resulting continuous sequence, the vascular term introduced a smooth baseline variation with an amplitude of up to $30\%$ of the peak pulsatile vascular change, while the muscle term comprised one to three weak, smoothly varying conductivity regions confined to the simulated muscle tissue, with a maximum amplitude of $0.006~\mathrm{S/m}$.
The combination of these terms prevented all temporal conductivity changes from being represented as identical scaled copies of the arterial pulse.

The first and seventh beats provided temporal padding for the 100-sample moving-mean decomposition.
For any continuous conductivity or voltage sequence $X_t$, let $\mathcal M_{100}(X)_t$ denote its 100-sample moving mean and let $t_{\mathrm{ref}}$ denote the first sample of the retained sequence.
The stored low-pass and high-pass audit components were
\begin{subequations}
\label{eq:gcnm-temporal-components}
\begin{align}
    X^{\mathrm{LP}}_t
    &=
    \mathcal M_{100}(X)_t-\mathcal M_{100}(X)_{t_{\mathrm{ref}}},
    \label{eq:gcnm-low-pass}\\
    X^{\mathrm{HP}}_t
    &=
    X_t-\mathcal M_{100}(X)_t
    -X_{t_{\mathrm{ref}}}
    +\mathcal M_{100}(X)_{t_{\mathrm{ref}}}.
    \label{eq:gcnm-high-pass}
\end{align}
\end{subequations}
Equations~\eqref{eq:gcnm-low-pass} and \eqref{eq:gcnm-high-pass} satisfy $X^{\mathrm{HP}}_t+X^{\mathrm{LP}}_t=X_t-X_{t_{\mathrm{ref}}}$.
The two padding beats were discarded after filtering and temporal decomposition, leaving five beats and 250 samples per anatomy.
This construction preserved complete pulse cycles while reducing boundary effects in the stored temporal components.

For every sample, the anatomical conductivity field was first evaluated on a fine forward mesh.
The corresponding eight-electrode voltages were then simulated using the skip-2 stimulation and skip-1 measurement patterns described in Supplementary Discussion~6.6.
A distinct coarser inverse mesh was used to store the target conductivity and to train the GCNM.
Using different forward and inverse meshes reduced the inverse crime that would arise if the same discretization were used to generate and reconstruct the measurements.
To reduce the computational cost of generating the training data, we computed one exact resting forward solution and one Jacobian for each virtual anatomy and approximated its time-varying voltages using
\begin{equation}
    V^{\mathrm{lin}}_{a,t}
    =
    F\!\left(\sigma^{\mathrm f}_{a,\mathrm{rest}}\right)
    +J_{a,\mathrm{rest}}
    \left(
        \sigma^{\mathrm f}_{a,t}
        -\sigma^{\mathrm f}_{a,\mathrm{rest}}
    \right).
\label{eq:gcnm-linearized-voltage}
\end{equation}
The first-order approximation in equation~\eqref{eq:gcnm-linearized-voltage} retained the heterogeneous resting conductivity of the corresponding virtual anatomy and replaced an exact nonlinear forward solution at every training sample.
Validation and test voltages were instead generated as $V^{\mathrm{nl}}_{a,t}=F(\sigma^{\mathrm f}_{a,t})$, allowing the linearized approximation to be evaluated on anatomies that were not used for training.
After applying the same filtering and temporal decomposition to both solutions, the approximation error for each component $c\in\{\mathrm{full},\mathrm{HP},\mathrm{LP}\}$ was quantified by
\begin{subequations}
\label{eq:gcnm-linearization-metrics}
\begin{align}
    \operatorname{NRMSE}_{a,c}
    &=
    \frac{
        \left\|
            V^{\mathrm{lin}}_{a,c}
            -V^{\mathrm{nl}}_{a,c}
        \right\|_2
    }{
        \left\|V^{\mathrm{nl}}_{a,c}\right\|_2
    },
    \label{eq:gcnm-linearization-nrmse}\\
    \rho_{a,c}
    &=
    \operatorname{corr}\!\left(
        \operatorname{vec}\!\left(V^{\mathrm{lin}}_{a,c}\right),
        \operatorname{vec}\!\left(V^{\mathrm{nl}}_{a,c}\right)
    \right).
    \label{eq:gcnm-linearization-correlation}
\end{align}
\end{subequations}
Using the metrics in equations~\eqref{eq:gcnm-linearization-nrmse} and \eqref{eq:gcnm-linearization-correlation}, a ring-specific corpus was accepted only when $\max_a\operatorname{NRMSE}_{a,c}\leq0.50$ and $\operatorname{mean}_a\rho_{a,c}\geq0.80$ for every component in both held-out partitions.

The simulator additionally generated channel gain, current gain, voltage offset, contact-drift, white-noise, correlated-noise, and transient-artifact variations.
Both the augmented and clean voltage sequences were filtered using the third-order zero-phase $5~\mathrm{Hz}$ low-pass filter used in the PVI processing pipeline, and both views were retained so that acquisition robustness could be studied without regenerating the anatomies.

Of the 19 ring meshes described in Supplementary Discussion~6.1, 15 were represented in the included cohort and were used for GCNM training.
These meshes ranged from US060 to US130 in half-size increments.
For each ring, 200 virtual anatomies were partitioned before beat and sample expansion into 160 training anatomies, 20 validation anatomies, and 20 test anatomies.
Each anatomy contributed five retained beats.
The resulting split contained 800 training beats, 100 validation beats, and 100 test beats per ring.
Since each beat contained 50 ordered samples, this yielded 40,000 training samples, 5,000 validation samples, and 5,000 test samples per ring.
No anatomy, beat, or sample was shared between the three partitions.


\clearpage
\color{black}

\paragraph{Ring-specific training}
\label{sec:gcnm-ring-training}
\mbox{}\par

One two-stage GCNM was trained for each of the 15 ring geometries.
For a given ring, the conductivity targets and voltages were paired with the forward mesh, inverse mesh, graph, and projected Laplacian associated with that geometry, ensuring that its electrode arrangement, finite-element discretization, physics-based update, and graph connectivity remained consistent throughout training and reconstruction.
The resulting 15 models contained 30 independently parameterized GCN blocks and each model was applied only to participants measured with its corresponding ring geometry.

Training followed the sequential procedure defined above.
For each ring, the first block was fitted using 40,000 samples and its checkpoint was selected on the independent 5,000-sample validation partition.
After freezing this checkpoint, we reconstructed the training and validation samples and recomputed the nonlinear forward solution, residual, Jacobian, and physics-based update from each first-stage output.
These recomputed quantities formed the inputs used to fit and select the second block, while the 5,000 exact nonlinear test samples remained excluded from parameter optimization and checkpoint selection.

Both blocks were optimized with Adam using an initial learning rate of $10^{-3}$ and batches of 512 graphs.
Training was limited to 150 epochs and stopped when the validation selection score did not improve for 30 consecutive epochs, with the random seed fixed at zero.
All ring-specific models used the same direct-output formulation, coordinate features, target and background weighting, and three hidden graph layers with 64 channels.

Before training all ring-specific models, we selected the architecture and loss coefficients in a controlled ring size 12.0 pilot shown in Supplementary Fig.~24.
This plot compared different GCN architectures and parameters across different metrics like the inclusion of element coordinates, target-magnitude weighting, and background weighting.
The selected formulation used the direct output with coordinates, $\alpha=1$, and $\beta=0.25$, while the Laplacian regularization parameter $\lambda=5\times10^{-4}$ and the 64-channel hidden width were held fixed rather than exhaustively optimized.

We compared the conventional Newton reconstruction and both GCNM stages with the known conductivity target on a standardized exact nonlinear visualization set containing three complete beats from each test anatomy for each ring.
Across the 15 ring sizes, the median NRMSE was $4.72$ for Newton, $1.12$ for the first GCNM stage, and $1.12$ for the second stage.
The corresponding median spatial correlations were $-0.013$, $-0.090$, and $-0.075$.
Visual inspection of Supplementary Fig.~25 showed that the GCNM stages recovered two spatially distinct regions corresponding to the synthetic arterial targets.
Relative to the first stage, the second stage provided a small aggregate improvement, with the mean NRMSE decreasing from $1.20$ to $1.19$ and the mean spatial correlation increasing from $-0.060$ to $-0.053$ across the 15 ring sizes.
Although this improvement was marginal, it indicates that recomputing the physics-based quantities after the first-stage reconstruction can provide additional refinement.
The usefulness of this second stage and of further reconstruction stages should be investigated in future work.

After reconstruction, the GCNM outputs were rasterized on the $40\times40$ grid defined in Supplementary Discussion~6.1.
Because direct conductivity ground truth was unavailable for the experimental recordings, their temporal behavior was compared qualitatively in the participant atlas shown in Supplementary Fig.~26.
The first complete 50-sample beat in a predefined evaluation segment for each participant will be used to avoid selecting examples according to their visual appearance.
Five approximately equally spaced samples with indices $1$, $10$, $20$, $30$, $40$ and $50$ will span the beat from its first to its final sample.
Spatial orientation, image extent, sample index, and conductivity units will be identical between the two rows.
A common color scale will be used across all patients and all five samples so that changes in amplitude and spatial concentration remain directly comparable.

\paragraph{Sequential training}
\label{sec:gcnm-sequential-training}

The two GCN blocks were trained sequentially rather than end-to-end.
After the first stage was trained, its output was used to recompute the nonlinear differential forward solution, residual, Jacobian, and physics-based update in equation~\eqref{eq:gcnm-physics-update} before training the second stage.
To define the implemented objective, let a mini-batch contain $Q$ graphs indexed by $q$, with the $N_K$ mesh elements in each graph indexed by $i$.
The normalized target and prediction were
\begin{equation*}
    y_{q,i}
    =
    \frac{\Delta\sigma^{*}_{q,i}}{s},
    \qquad
    \widehat y_{q,i}
    =
    \frac{1}{s}
    \left[
        \operatorname{GCN}_{\theta_n}\!\left(H_{n,q},G\right)
    \right]_i.
\end{equation*}
Here, $y_{q,i}$ is the normalized target differential conductivity for element $i$ in graph $q$, $\widehat y_{q,i}$ is the corresponding normalized GCN prediction, $\Delta\sigma^{*}_{q,i}$ is the target differential conductivity before normalization, and $s$ is the conductivity scale defined above.
For each graph, the reference target magnitude was
\begin{equation*}
    m_q
    =
    \underset{j:\,|y_{q,j}|>0}{\operatorname{median}}
    |y_{q,j}|.
\end{equation*}
Here, $m_q$ is the median magnitude of the nonzero normalized target elements in graph $q$.
If a graph contained no nonzero target element, the implementation set $m_q=1$.
The elementwise weight was
\begin{equation*}
    w_{q,i}
    =
    1+\alpha\,
    \operatorname{clip}\!\left(
        \frac{|y_{q,i}|}{\max(m_q,\varepsilon)},\,0,\,2
    \right),
    \qquad
    \varepsilon=10^{-8}.
\end{equation*}
Here, $w_{q,i}$ is the target-dependent weight for element $i$, $\alpha$ controls the additional weight assigned to nonzero and higher-magnitude target elements, and $\varepsilon$ is the numerical threshold used to identify target-background elements.
The target-background set for graph $q$ was
\begin{equation*}
    \mathcal B_q
    =
    \left\{i:\ |y_{q,i}|\leq\varepsilon\right\}.
\end{equation*}
Here, $\mathcal B_q$ contains the mesh elements whose normalized target magnitude did not exceed $\varepsilon$.
The loss minimized at each stage was
\begingroup
\small
\begin{equation}
    \mathcal L_n(\theta_n)
    =
    \frac{1}{QN_K}\sum_{q=1}^{Q}\sum_{i=1}^{N_K}
    w_{q,i}\left(\widehat y_{q,i}-y_{q,i}\right)^2
    +
    \frac{\beta}{\sum_{q=1}^{Q}|\mathcal B_q|}
    \sum_{q=1}^{Q}\sum_{i\in\mathcal B_q}
    \left(\widehat y_{q,i}-y_{q,i}\right)^2.
\label{eq:gcnm-training-loss}
\end{equation}
\endgroup
\begin{equation*}
    \alpha=1,
    \qquad
    \beta=0.25.
\end{equation*}
The first term is the target-weighted mean squared error over all $QN_K$ mesh elements.
The second term is the mean squared error over the target-background sets, and $\beta$ controls its contribution to the total loss.
The loss in equation~\eqref{eq:gcnm-training-loss} and the same coefficients were used for both stages.
Both GCNs were optimized with Adam (initial learning rate $10^{-3}$), and their parameters were selected independently.
The saved validation checkpoint minimized the composite score
\begin{equation}
    S_{\mathrm{val}}
    =
    \operatorname{NRMSE}_{\mathrm{val}}
    +
    \operatorname{RMS}_{\mathrm{background,val}}
    -
    0.25\,\operatorname{Dice}_{\mathrm{support,val}}.
\label{eq:gcnm-validation-score}
\end{equation}
To define the three terms in equation~\eqref{eq:gcnm-validation-score} explicitly, let $Q_{\mathrm{val}}$ be the number of graphs in the validation partition and let $\widehat y_{q,i}$ denote the normalized prediction from the checkpoint being evaluated.
For each validation graph, the target support, its cardinality, and the equal-cardinality predicted support were
\begin{equation*}
    \mathcal S_q
    =
    \left\{i:\ |y_{q,i}|>\varepsilon\right\}
    =
    \left\{1,\ldots,N_K\right\}\setminus\mathcal B_q,
    \qquad
    K_q=|\mathcal S_q|,
    \qquad
    \widehat{\mathcal S}_q
    =
    \operatorname{TopK}_{K_q}
    \left(\left\{|\widehat y_{q,i}|\right\}_{i=1}^{N_K}\right),
\end{equation*}
where $\operatorname{TopK}_{K_q}$ returns the indices of the $K_q$ largest absolute predicted values.
Let $\mathcal V_{+}=\{q:\ K_q>0\}$ denote the validation graphs having nonempty target support.
The implemented validation metrics were
\begin{subequations}
\label{eq:gcnm-validation-metrics}
\begin{align}
    \operatorname{NRMSE}_{\mathrm{val}}
    &=
    \frac{
        \left[
            \dfrac{1}{Q_{\mathrm{val}}N_K}
            \sum_{q=1}^{Q_{\mathrm{val}}}
            \sum_{i=1}^{N_K}
            \left(\widehat y_{q,i}-y_{q,i}\right)^2
        \right]^{1/2}
    }{
        \max\!\left(
            \left[
                \dfrac{1}{Q_{\mathrm{val}}N_K}
                \sum_{q=1}^{Q_{\mathrm{val}}}
                \sum_{i=1}^{N_K}
                y_{q,i}^{\,2}
            \right]^{1/2},
            10^{-12}
        \right)
    },
    \label{eq:gcnm-validation-nrmse}\\
    \operatorname{RMS}_{\mathrm{background,val}}
    &=
    \left[
        \frac{
            \displaystyle
            \sum_{q=1}^{Q_{\mathrm{val}}}
            \sum_{i\in\mathcal B_q}
            \left(\widehat y_{q,i}-y_{q,i}\right)^2
        }{
            \displaystyle
            \sum_{q=1}^{Q_{\mathrm{val}}}|\mathcal B_q|
        }
    \right]^{1/2},
    \label{eq:gcnm-validation-background-rms}\\
    \operatorname{Dice}_{\mathrm{support,val}}
    &=
    \frac{1}{|\mathcal V_{+}|}
    \sum_{q\in\mathcal V_{+}}
    \frac{
        2\left|\mathcal S_q\cap\widehat{\mathcal S}_q\right|
    }{
        |\mathcal S_q|+|\widehat{\mathcal S}_q|
    }
    =
    \frac{1}{|\mathcal V_{+}|}
    \sum_{q\in\mathcal V_{+}}
    \frac{
        \left|\mathcal S_q\cap\widehat{\mathcal S}_q\right|
    }{
        K_q
    }.
    \label{eq:gcnm-validation-support-dice}
\end{align}
\end{subequations}
The NRMSE and background RMS were calculated after concatenating all validation elements, whereas the support Dice was calculated separately for each graph and then averaged over graphs with nonempty support.
The final equality in equation~\eqref{eq:gcnm-validation-support-dice} follows because $|\widehat{\mathcal S}_q|=|\mathcal S_q|=K_q$.
All quantities in equations~\eqref{eq:gcnm-validation-metrics} used the normalized targets and predictions; consequently, the background RMS was dimensionless and could be combined directly with the NRMSE and Dice in equation~\eqref{eq:gcnm-validation-score}.
The validation score in equation~\eqref{eq:gcnm-validation-score} was used only for checkpoint selection and was not an additional differentiable training-loss term.


\clearpage
\begingroup
\color{black}

\begin{figure}[p]
    \centering
    \textbf{Supplementary Fig.~24. Controlled ring-size-12 architecture and loss ablation.}\par
    \smallskip
    \begin{minipage}{0.94\textwidth}
        \small
        Comparison between different GCNM model configurations using multiple performance metrics.
        Blue identifies the GCNM input or output formulation, including direct, residual, and
        coordinate-augmented configurations. Orange identifies the value of the $\alpha$ target-weighting
        coefficient, and green identifies the background-weighting coefficient $\beta$. The selected
        configuration used the direct coordinate formulation with $\alpha=1$ and $\beta=0.25$.
    \end{minipage}\par
    \vspace{0.8em}
    \includegraphics[
        width=0.98\textwidth,
        height=0.70\textheight,
        keepaspectratio
    ]{figures/fig2_final.png}
\end{figure}

\clearpage
\begin{figure}[p]
    \centering
    \textbf{Supplementary Fig.~25. Ground-truth simulated graph-convolutional-network finger-image reconstructions.}\par
    \smallskip
    \begin{minipage}{0.94\textwidth}
        \small
        Synthetic examples from ring meshes of sizes 6 through 13, each with a different virtual finger
        anatomy, are shown with the corresponding ground-truth and second-stage GCNM conductivity
        reconstructions at time instants $t_i$ within one beat, with $i=1,2,\ldots,50$. the time sample within a beat period. 
    \end{minipage}\par
    \vspace{0.8em}
    \includegraphics[
        width=0.98\textwidth,
        height=0.70\textheight,
        keepaspectratio
    ]{figures/fig1_final.png}
\end{figure}

\clearpage
\begin{figure}[p]
    \centering
    \textbf{Supplementary Fig.~26. Finger conductivity image reconstruction with a graph convolutional network.}\par
    \smallskip
    \begin{minipage}{0.94\textwidth}
        \small
        Participant-level temporal progression of second-stage GCNM reconstructions
        at six samples spanning one complete beat. The same spatial orientation, image extent, sample
        indices, conductivity units, and participant-specific color scale are used for both methods.
    \end{minipage}\par
    \vspace{0.8em}
    \includegraphics[
        width=0.90\textwidth,
        height=0.72\textheight,
        keepaspectratio
    ]{figures/fig3_final.png}
\end{figure}

\endgroup
\clearpage

\endgroup

\endgroup

\begin{thebibliography}{9}

\bibitem{Herzberg2021GCNM}
W.~Herzberg, D.~B. Rowe, A.~Hauptmann, and S.~J. Hamilton, ``Graph convolutional networks for model-based learning in nonlinear inverse problems,'' \emph{IEEE Transactions on Computational Imaging}, vol.~7, pp.~1341\textendash 1353, 2021.

\bibitem{ITIS2024}
C.~Baumgartner, P.~A. Hasgall, F.~Di Gennaro, E.~Neufeld, B.~Lloyd, M.~C. Gosselin, D.~Payne, A.~Klingenb\"ock, and N.~Kuster, ``IT'IS Database for Thermal and Electromagnetic Parameters of Biological Tissues,'' version~4.2, IT'IS Foundation, 2024, doi: 10.13099/VIP21000-04-2.

\end{thebibliography}

\end{document}

